\section{Movimiento circular}

Hasta ahora hemos estudiado movimientos rectilíneos y movimientos curvilíneos,
como el de los proyectiles. En estos últimos, la trayectoria no es una línea
recta y, por tanto, la dirección del vector velocidad puede cambiar
continuamente.

Existe un tipo particular de movimiento curvilíneo en el que la trayectoria es
una circunferencia: el \textbf{movimiento circular}.

Podemos observarlo en numerosos fenómenos y dispositivos:

\begin{itemize}
    \item las aspas de un ventilador;
    \item los puntos de una rueda que gira alrededor de un eje fijo;
    \item el extremo de la manecilla de un reloj;
    \item una piedra que gira sujeta a una cuerda;
    \item las piezas de una centrifugadora;
    \item una cabina de una rueda de observación.
\end{itemize}

Aunque estos movimientos tienen características diferentes, comparten una
propiedad geométrica: cada partícula considerada describe una trayectoria
circular alrededor de un centro.

Surgen entonces algunas preguntas importantes.

¿Por qué un objeto puede acelerar aunque su rapidez permanezca constante?
¿Cómo se relaciona el número de vueltas con la velocidad? ¿Qué diferencia
existe entre la velocidad de un punto cercano al centro de una rueda y la de
otro situado cerca del borde?

Para responderlas necesitamos complementar las magnitudes lineales ya conocidas
con nuevas magnitudes que describan el cambio de posición angular.


\begin{idea}[Idea fundamental]

El movimiento circular puede describirse mediante magnitudes lineales y
angulares.

Las magnitudes lineales permiten estudiar la posición, la velocidad y la
aceleración de la partícula en el espacio.

Las magnitudes angulares permiten describir cómo cambia su posición alrededor
del centro de la circunferencia.

Ambas descripciones corresponden al mismo movimiento y deben ser coherentes
entre sí.

\end{idea}


\section{Concepto de movimiento circular}

\begin{definicion}[Movimiento circular]

El movimiento circular es aquel en el que una partícula describe una
trayectoria circular respecto de un sistema de referencia determinado.

Durante este movimiento, la distancia entre la partícula y el centro de la
circunferencia permanece constante.

\end{definicion}

Si llamamos \(R\) al radio de la trayectoria y \(O\) a su centro, entonces:

\[
OP=R=\text{constante}.
\]

Esta condición geométrica distingue al movimiento circular de otros
movimientos curvilíneos.

Sin embargo, no implica que la rapidez de la partícula sea constante.

Una partícula puede recorrer una circunferencia:

\begin{itemize}
    \item con rapidez constante;
    \item aumentando su rapidez;
    \item disminuyendo su rapidez;
    \item cambiando de sentido.
\end{itemize}

Por tanto, debemos distinguir entre la forma de la trayectoria y la manera en
que la partícula se desplaza sobre ella.


\section{Modelo físico y sistema de referencia}

Consideraremos una partícula que se mueve sobre una circunferencia fija de
radio \(R\).

\begin{modelo}[Partícula en trayectoria circular]

Supondremos que:

\begin{itemize}
    \item el objeto puede representarse mediante una partícula;
    \item el centro de la circunferencia permanece fijo;
    \item el radio de la trayectoria es constante;
    \item el movimiento ocurre en un plano fijo;
    \item las magnitudes cinemáticas pueden variar con el tiempo.
\end{itemize}

\end{modelo}

Establecemos un sistema cartesiano \(xy\):

\begin{itemize}
    \item el origen coincide con el centro \(O\);
    \item el eje \(+x\) apunta hacia la derecha;
    \item el eje \(+y\) apunta hacia arriba;
    \item los vectores unitarios son \(\ihat\) y \(\jhat\);
    \item los ángulos positivos se miden en sentido antihorario.
\end{itemize}

\begin{center}
\begin{tikzpicture}[scale=1.0]

    \coordinate (O) at (0,0);
    \coordinate (P) at (2.4,1.8);

    \draw[-{Latex},thick]
        (-3.5,0) -- (3.8,0)
        node[right] {\(x\)};

    \draw[-{Latex},thick]
        (0,-3.3) -- (0,3.5)
        node[above] {\(y\)};

    \draw[
        very thick,
        AzulPrincipal
    ]
        (O) circle (3);

    \draw[-{Latex},very thick,VerdePrincipal]
        (O) -- (P)
        node[midway,above left] {\(\vect{r}\)};

    \draw[dashed]
        (P) -- (2.4,0);

    \draw[dashed]
        (P) -- (0,1.8);

    \fill (O) circle (2pt);
    \node[below left] at (O) {\(O\)};

    \fill (P) circle (3pt);
    \node[above right] at (P) {\(P\)};

    \node[below] at (2.4,0) {\(x\)};
    \node[left] at (0,1.8) {\(y\)};

    \draw[
        -{Latex},
        thick,
        RojoPrincipal
    ]
        (0.85,0)
        arc[
            start angle=0,
            end angle=36.87,
            radius=0.85
        ];

    \node[RojoPrincipal] at (1.15,0.38) {\(\theta\)};

    \node at (-2.1,2.5) {\(R\)};

\end{tikzpicture}
\end{center}

El radio \(R\) es una magnitud escalar positiva:

\[
R>0.
\]

El vector de posición cambia de dirección mientras la partícula gira, aunque
su magnitud permanezca constante:

\[
\norm{\vect{r}}=R.
\]

\begin{interpretacion}

La partícula no mantiene constante su vector de posición.

Lo que permanece constante es únicamente su magnitud.

Por ello, incluso cuando el radio no cambia, existe movimiento porque cambia
la dirección de \(\vect{r}\).

\end{interpretacion}


\section{Posición angular}

En una trayectoria rectilínea la posición puede describirse mediante una
coordenada \(x\).

En una trayectoria circular resulta especialmente conveniente utilizar un
ángulo.

\begin{definicion}[Posición angular]

La posición angular \(\theta\) es el ángulo orientado que forma el radio
dirigido hacia la partícula con una dirección de referencia previamente
establecida.

\end{definicion}

Tomaremos como dirección de referencia el eje \(+x\).

La posición angular se mide positivamente en sentido antihorario.

Por tanto:

\[
\theta>0
\]

representa un giro antihorario desde la dirección de referencia, mientras que
un giro horario corresponde a una variación angular negativa.

La posición angular puede describirse mediante una función:

\[
\boxed{\theta=\theta(t)}.
\]

\begin{advertencia}

La posición angular no determina cuántas vueltas ha completado una partícula si
solamente se informa su orientación final dentro de una vuelta.

Por ejemplo, una partícula situada nuevamente sobre el eje \(+x\) podría
haber completado cero, una, dos o muchas vueltas.

Para describir el movimiento acumulado necesitamos registrar el ángulo sin
reducirlo automáticamente al intervalo de una revolución.

\end{advertencia}


\section{El radián}

Para relacionar magnitudes angulares con longitudes necesitamos una unidad
angular especialmente útil: el radián.

Consideremos una circunferencia de radio \(R\).

Un ángulo central determina sobre ella un arco de longitud \(s\).

\begin{definicion}[Radián]

Un radián es el ángulo central que subtiende un arco cuya longitud es igual al
radio de la circunferencia.

\end{definicion}

Si:

\[
s=R,
\]

entonces:

\[
\theta=\SI{1}{\radian}.
\]

En general:

\[
\boxed{\theta=\frac{s}{R}}.
\]

Esta relación constituye la definición de la medida angular en radianes.

Despejando:

\[
\boxed{s=R\theta}.
\]

\begin{condicion}

La relación:

\[
s=R\theta
\]

requiere que el ángulo esté expresado en radianes.

Además, \(s\) representa una longitud de arco medida con el signo o la
orientación apropiados cuando se utilice como desplazamiento angular.

\end{condicion}


\subsection{Relación entre radianes, grados y revoluciones}

Una revolución completa corresponde a:

\[
360\grado.
\]

Como la longitud de una circunferencia es:

\[
2\pi R,
\]

el ángulo correspondiente es:

\[
\theta
=
\frac{2\pi R}{R}
=
2\pi.
\]

Por tanto:

\[
\boxed{
1\text{ revolución}
=
2\pi\ \si{\radian}
=
360\grado
}.
\]

De aquí:

\[
\boxed{
180\grado=\pi\ \si{\radian}
}.
\]

Las conversiones son:

\[
\theta_{\mathrm{rad}}
=
\theta_{\mathrm{grados}}
\frac{\pi}{180},
\]

\[
\theta_{\mathrm{grados}}
=
\theta_{\mathrm{rad}}
\frac{180}{\pi}.
\]

\begin{interpretacion}

El radián es una unidad coherente del SI para expresar ángulos planos.

Aunque dimensionalmente el cociente \(s/R\) es adimensional, escribir
\(\si{\radian}\) permite distinguir una magnitud angular de otras cantidades
adimensionales.

\end{interpretacion}


\section{Desplazamiento angular}

Cuando una partícula pasa de una posición angular inicial:

\[
\theta_i
\]

a una posición angular final:

\[
\theta_f,
\]

definimos su desplazamiento angular.

\begin{definicion}[Desplazamiento angular]

El desplazamiento angular es la diferencia entre la posición angular final y
la inicial:

\[
\boxed{
\Delta\theta=\theta_f-\theta_i
}.
\]

\end{definicion}

Si:

\[
\Delta\theta>0,
\]

el desplazamiento angular neto corresponde al sentido antihorario.

Si:

\[
\Delta\theta<0,
\]

corresponde al sentido horario.

En un movimiento que puede completar varias vueltas, utilizaremos una
posición angular continua que registre los giros acumulados.


\subsection{Longitud de arco y desplazamiento angular}

Si el radio permanece constante, la longitud de arco orientada asociada al
desplazamiento angular es:

\[
\boxed{\Delta s=R\Delta\theta}.
\]

Si el movimiento no cambia de sentido durante el intervalo, la distancia
recorrida sobre la circunferencia es:

\[
\boxed{d=R|\Delta\theta|}.
\]

Si existen cambios de sentido, debemos sumar las longitudes de todos los arcos
recorridos.

\begin{errorfrecuente}[Arco y desplazamiento rectilíneo]

La longitud del arco recorrido no es igual, en general, a la magnitud del
vector desplazamiento:

\[
\norm{\Delta\vect{r}}.
\]

El arco sigue la circunferencia.

El vector desplazamiento une directamente las posiciones inicial y final.

\end{errorfrecuente}


\section{Distancia recorrida y desplazamiento en una vuelta}

Supongamos que una partícula completa una revolución y regresa exactamente a
su posición inicial.

La distancia recorrida es la longitud de la circunferencia:

\[
\boxed{d=2\pi R}.
\]

El desplazamiento angular acumulado, para una vuelta antihoraria, es:

\[
\boxed{\Delta\theta=2\pi\ \si{\radian}}.
\]

Pero el vector desplazamiento es:

\[
\Delta\vect{r}
=
\vect{r}_f-\vect{r}_i.
\]

Como:

\[
\vect{r}_f=\vect{r}_i,
\]

obtenemos:

\[
\boxed{\Delta\vect{r}=\vect{0}}.
\]

\begin{interpretacion}

Después de una vuelta completa:

\[
d=2\pi R,
\]

pero:

\[
\norm{\Delta\vect{r}}=0.
\]

Esto demuestra que la distancia recorrida, el desplazamiento vectorial y el
desplazamiento angular representan magnitudes diferentes.

\end{interpretacion}


\section{Ejemplo: desplazamiento angular}

\begin{ejemplo}

Una rueda de radio:

\[
R=\SI{0.40}{\meter}
\]

gira un ángulo de:

\[
150\grado
\]

en sentido antihorario.

Determine:

\begin{enumerate}[label=\alph*)]
    \item el desplazamiento angular en radianes;
    \item la longitud del arco recorrido por un punto del borde.
\end{enumerate}

\end{ejemplo}

Primero convertimos el ángulo:

\[
\Delta\theta
=
150\frac{\pi}{180}.
\]

Entonces:

\[
\boxed{
\Delta\theta
=
\frac{5\pi}{6}\ \si{\radian}
\approx
\SI{2.62}{\radian}
}.
\]

Como el punto se mueve sobre una circunferencia de radio constante:

\[
s=R\Delta\theta.
\]

Por tanto:

\[
s
=
(\SI{0.40}{\meter})
\left(
\frac{5\pi}{6}
\right).
\]

Así:

\[
\boxed{s\approx\SI{1.05}{\meter}}.
\]

Este resultado representa la longitud del arco recorrido, no la distancia
recta entre los extremos.


\section{Velocidad angular media}

Una vez descrita la posición angular podemos preguntar qué tan rápidamente
cambia con el tiempo.

\begin{definicion}[Velocidad angular media]

La velocidad angular media es el desplazamiento angular dividido entre el
intervalo de tiempo:

\[
\boxed{
\bar{\omega}
=
\frac{\Delta\theta}{\Delta t}
}.
\]

\end{definicion}

La unidad SI es:

\[
\si{\radian\per\second}.
\]

La velocidad angular media puede ser positiva o negativa.

En nuestra convención:

\[
\bar{\omega}>0
\]

indica un desplazamiento angular neto antihorario.

En cambio:

\[
\bar{\omega}<0
\]

indica un desplazamiento angular neto horario.

\begin{interpretacion}

Si:

\[
\bar{\omega}
=
\SI{2}{\radian\per\second},
\]

significa que el desplazamiento angular aumenta, en promedio, dos radianes por
cada segundo transcurrido.

No significa necesariamente que la partícula mantenga una velocidad angular
constante durante todo el intervalo.

\end{interpretacion}


\section{Velocidad angular instantánea}

La velocidad angular media describe un intervalo temporal.

Para conocer el cambio angular en un instante, tomamos el límite cuando el
intervalo se hace muy pequeño.

\begin{definicion}[Velocidad angular instantánea]

La velocidad angular instantánea es:

\[
\boxed{
\omega
=
\lim_{\Delta t\to0}
\frac{\Delta\theta}{\Delta t}
=
\deriv{\theta}{t}
}.
\]

\end{definicion}

Esta derivada expresa la tasa instantánea de cambio de la posición angular.

En movimiento circular plano utilizaremos \(\omega\) como cantidad algebraica
con signo, respecto del eje de rotación previamente definido.

\begin{advertencia}

No debe confundirse:

\[
\omega
\]

con:

\[
\vect{\omega}.
\]

La primera representa aquí la componente angular con signo.

El vector velocidad angular, cuando se utilice, apunta a lo largo del eje de
rotación y su sentido se determina mediante la regla de la mano derecha.

\end{advertencia}


\section{Velocidad angular vectorial}

Para una trayectoria circular situada en el plano \(xy\), el eje de rotación
es perpendicular a dicho plano.

Definimos:

\[
\boxed{
\vect{\omega}
=
\omega\khat
}.
\]

Si el giro es antihorario visto desde \(+z\):

\[
\omega>0,
\]

y el vector velocidad angular apunta hacia:

\[
+\khat.
\]

Si el giro es horario:

\[
\omega<0,
\]

y apunta hacia:

\[
-\khat.
\]

La magnitud del vector es:

\[
\norm{\vect{\omega}}=|\omega|.
\]

\begin{interpretacion}

El vector velocidad angular no apunta tangente a la trayectoria.

Su dirección es perpendicular al plano del movimiento.

Esta diferencia es fundamental:

\[
\vect{v}
\]

es tangente a la circunferencia, mientras:

\[
\vect{\omega}
\]

apunta a lo largo del eje de rotación.

\end{interpretacion}


\section{Relación entre velocidad lineal y angular}

Consideremos una partícula que recorre un arco \(s\) de una circunferencia de
radio constante \(R\).

Sabemos que:

\[
s=R\theta.
\]

Para describir el movimiento con signo podemos escribir:

\[
\dd s=R\,\dd\theta,
\]

donde \(\dd s\) representa el incremento orientado de arco en el sentido
angular positivo.

Dividiendo entre \(\dd t\):

\[
\deriv{s}{t}
=
R\deriv{\theta}{t}.
\]

Como:

\[
\omega=\deriv{\theta}{t},
\]

obtenemos:

\[
\boxed{
v_{\mathrm{tan}}=R\omega
}.
\]

Aquí \(v_{\mathrm{tan}}\) es la componente de velocidad respecto del sentido
tangencial positivo.

La rapidez, que no puede ser negativa, es:

\[
\boxed{
v=R|\omega|
}.
\]

\begin{interpretacion}

Para una misma rapidez angular, los puntos situados a mayor distancia del
centro poseen mayor rapidez lineal.

Por ejemplo, en una rueda rígida que gira alrededor de un eje fijo:

\begin{itemize}
    \item todos los puntos completan una vuelta en el mismo tiempo;
    \item todos poseen la misma velocidad angular;
    \item los puntos cercanos al borde recorren arcos más largos;
    \item por ello, poseen mayor rapidez lineal.
\end{itemize}

\end{interpretacion}


\section{Relación vectorial entre velocidad angular y velocidad lineal}

La relación anterior también puede escribirse vectorialmente:

\[
\boxed{
\vect{v}
=
\vect{\omega}\pcross\vect{r}
}.
\]

Esta expresión es válida para una partícula que describe una circunferencia
alrededor del origen fijo, donde \(\vect{r}\) se mide desde el centro.

Como:

\[
\vect{\omega}
=
\omega\khat,
\]

y:

\[
\vect{r}
=
x\ihat+y\jhat,
\]

tenemos:

\[
\vect{v}
=
\omega\khat
\pcross
(x\ihat+y\jhat).
\]

Utilizando:

\[
\khat\pcross\ihat=\jhat,
\]

\[
\khat\pcross\jhat=-\ihat,
\]

obtenemos:

\[
\boxed{
\vect{v}
=
-\omega y\ihat+\omega x\jhat
}.
\]

Por tanto, las componentes rectangulares son:

\[
\boxed{
v_x=-\omega y
},
\]

\[
\boxed{
v_y=\omega x
}.
\]

Estas ecuaciones muestran que la velocidad es perpendicular al vector radial.


\section{Posición cartesiana de una partícula en movimiento circular}

Volvamos a la geometría de la circunferencia.

Si la partícula se encuentra a una posición angular \(\theta\), entonces:

\[
x=R\cos\theta,
\]

\[
y=R\sen\theta.
\]

Por tanto:

\[
\boxed{
\vect{r}(t)
=
R\cos\theta(t)\,\ihat
+
R\sen\theta(t)\,\jhat
}.
\]

Esta ecuación permite expresar el movimiento circular en componentes
rectangulares.

Su importancia es que conserva explícitamente la dependencia temporal de la
posición.

\begin{interpretacion}

La partícula se mueve en dos dimensiones, aunque su posición pueda describirse
mediante una sola coordenada angular.

Esto ocurre porque la restricción:

\[
x^2+y^2=R^2
\]

relaciona las dos coordenadas cartesianas.

\end{interpretacion}


\section{Deducción de las componentes de velocidad}

Partimos de:

\[
x(t)=R\cos\theta(t),
\]

\[
y(t)=R\sen\theta(t).
\]

La velocidad instantánea se obtiene derivando respecto del tiempo.

Para la componente horizontal:

\[
v_x=\deriv{x}{t}.
\]

Aplicando la regla de la cadena:

\[
v_x
=
-R\sen\theta
\deriv{\theta}{t}.
\]

Como:

\[
\deriv{\theta}{t}=\omega,
\]

obtenemos:

\[
\boxed{
v_x=-R\omega\sen\theta
}.
\]

Para la componente vertical:

\[
v_y=\deriv{y}{t}.
\]

Entonces:

\[
\boxed{
v_y=R\omega\cos\theta
}.
\]

El vector velocidad es:

\[
\boxed{
\vect{v}
=
-R\omega\sen\theta\,\ihat
+
R\omega\cos\theta\,\jhat
}.
\]

Calculamos su magnitud:

\[
v
=
\sqrt{v_x^2+v_y^2}.
\]

Sustituyendo:

\[
v
=
\sqrt{
R^2\omega^2\sen^2\theta
+
R^2\omega^2\cos^2\theta
}.
\]

Factorizando:

\[
v
=
R|\omega|
\sqrt{\sen^2\theta+\cos^2\theta}.
\]

Como:

\[
\sen^2\theta+\cos^2\theta=1,
\]

resulta:

\[
\boxed{
v=R|\omega|
}.
\]

Hemos recuperado la relación entre rapidez lineal y rapidez angular mediante
componentes rectangulares.


\section{Período y frecuencia}

Muchos movimientos circulares se repiten.

Por ejemplo, una rueda que gira uniformemente pasa por una misma posición
angular una vez en cada vuelta.

Para describir esta repetición introducimos el período y la frecuencia.


\subsection{Período}

\begin{definicion}[Período]

El período \(T\) es el tiempo necesario para completar una revolución en un
movimiento circular periódico.

\end{definicion}

Su unidad SI es:

\[
\si{\second}.
\]

Si una partícula completa \(N\) vueltas en un tiempo \(\Delta t\), con vueltas
de igual duración:

\[
\boxed{
T=\frac{\Delta t}{N}
}.
\]


\subsection{Frecuencia}

\begin{definicion}[Frecuencia]

La frecuencia \(f\) es el número de ciclos completos realizados por unidad de
tiempo en un movimiento periódico.

\end{definicion}

Su unidad SI es el hercio:

\[
\SI{1}{\hertz}
=
\SI{1}{\per\second}.
\]

Para un movimiento periódico:

\[
\boxed{
f=\frac{1}{T}
}.
\]

También:

\[
\boxed{
T=\frac{1}{f}
}.
\]

\begin{interpretacion}

El período indica cuánto dura una vuelta.

La frecuencia indica cuántas vueltas se completan por segundo.

Son magnitudes inversamente proporcionales.

\end{interpretacion}


\section{Relación entre velocidad angular, período y frecuencia}

En una revolución completa:

\[
|\Delta\theta|=2\pi\ \si{\radian}.
\]

Si el movimiento circular es uniforme, el tiempo correspondiente es \(T\).

Por tanto:

\[
|\omega|
=
\frac{2\pi}{T}.
\]

Como:

\[
f=\frac{1}{T},
\]

obtenemos:

\[
\boxed{
|\omega|=2\pi f
}.
\]

Para movimiento antihorario, con nuestra convención:

\[
\omega=2\pi f.
\]

Para movimiento horario:

\[
\omega=-2\pi f.
\]

La rapidez lineal será:

\[
v=R|\omega|.
\]

Por tanto:

\[
\boxed{
v=\frac{2\pi R}{T}
}.
\]

También:

\[
\boxed{
v=2\pi Rf
}.
\]

Estas relaciones son válidas para movimiento circular uniforme.

\begin{advertencia}

Si la velocidad angular cambia continuamente, las expresiones:

\[
|\omega|=\frac{2\pi}{T},
\qquad
|\omega|=2\pi f
\]

no deben utilizarse indiscriminadamente con un supuesto período constante.

Un movimiento circular no uniforme puede no ser periódico.

\end{advertencia}


\section{Revoluciones por minuto}

En numerosos dispositivos se expresa la rapidez de giro mediante
revoluciones por minuto:

\[
\mathrm{rpm}.
\]

Esta unidad indica cuántas revoluciones se completan en un minuto.

Como:

\[
1\text{ revolución}=2\pi\ \si{\radian},
\]

y:

\[
1\text{ minuto}=\SI{60}{\second},
\]

tenemos:

\[
\boxed{
\SI{1}{rpm}
=
\frac{2\pi}{60}
\ \si{\radian\per\second}
}.
\]

Por tanto, si \(n\) representa el valor numérico de las revoluciones por
minuto:

\[
\boxed{
|\omega|
=
\frac{2\pi n}{60}
}.
\]

El resultado se expresa en:

\[
\si{\radian\per\second}.
\]

La dirección del giro debe especificarse por separado.


\section{Ejemplo: ventilador en movimiento circular uniforme}

\begin{ejemplo}

Un ventilador gira a:

\[
120\ \mathrm{rpm}.
\]

Un punto de una de sus aspas se encuentra a:

\[
R=\SI{0.15}{\meter}
\]

del eje de rotación.

Determine:

\begin{enumerate}[label=\alph*)]
    \item la frecuencia;
    \item el período;
    \item la rapidez angular;
    \item la rapidez lineal del punto.
\end{enumerate}

\end{ejemplo}

Primero convertimos las revoluciones por minuto:

\[
f
=
\frac{120}{60}
\ \si{\per\second}.
\]

Entonces:

\[
\boxed{f=\SI{2.0}{\hertz}}.
\]

El período es:

\[
T=\frac{1}{f}.
\]

Por tanto:

\[
\boxed{T=\SI{0.50}{\second}}.
\]

La rapidez angular es:

\[
|\omega|=2\pi f.
\]

Entonces:

\[
|\omega|
=
2\pi(2.0).
\]

Así:

\[
\boxed{
|\omega|
\approx
\SI{12.6}{\radian\per\second}
}.
\]

La rapidez lineal del punto es:

\[
v=R|\omega|.
\]

Por tanto:

\[
v
=
(\SI{0.15}{\meter})
(\SI{12.6}{\radian\per\second}).
\]

Así:

\[
\boxed{
v\approx
\SI{1.88}{\meter\per\second}
}.
\]

\begin{interpretacion}

El ventilador completa dos vueltas por segundo.

Todos los puntos de las aspas que giran rígidamente alrededor del mismo eje
poseen la misma rapidez angular.

Sin embargo, un punto más próximo al eje tiene menor rapidez lineal.

\end{interpretacion}


\section{Movimiento circular uniforme}

Ahora podemos definir uno de los modelos principales de este capítulo.

\begin{definicion}[Movimiento circular uniforme]

El movimiento circular uniforme, abreviado MCU, es aquel en el que una
partícula describe una trayectoria circular con rapidez constante.

\end{definicion}

En este movimiento:

\[
R=\text{constante},
\]

\[
v=\text{constante},
\]

\[
|\omega|=\text{constante}.
\]

Mientras el movimiento conserva su sentido, la velocidad angular algebraica
también permanece constante:

\[
\omega=\text{constante}.
\]

Sin embargo, debemos distinguir dos afirmaciones:

\[
\text{rapidez constante}
\]

y:

\[
\text{velocidad vectorial constante}.
\]

No significan lo mismo.


\section{¿Existe aceleración en el MCU?}

Supongamos que una partícula gira uniformemente alrededor de una circunferencia.

En diferentes puntos de la trayectoria, los vectores velocidad tienen igual
magnitud:

\[
v_1=v_2=v.
\]

Pero sus direcciones son diferentes:

\[
\vect{v}_1\neq\vect{v}_2.
\]

Por tanto:

\[
\Delta\vect{v}
=
\vect{v}_2-\vect{v}_1
\neq\vect{0}.
\]

Como existe un cambio del vector velocidad, existe aceleración.

\begin{principio}[Aceleración en el movimiento circular]

Toda partícula que se mueve siguiendo una trayectoria circular con rapidez
distinta de cero experimenta una aceleración asociada al cambio de dirección de
su velocidad.

En el movimiento circular uniforme, dicha aceleración apunta hacia el centro
de la circunferencia.

\end{principio}

\begin{idea}

Movimiento uniforme no significa necesariamente aceleración nula.

En movimiento rectilíneo uniforme, la velocidad vectorial permanece constante.

En movimiento circular uniforme, sólo permanece constante su magnitud.

\end{idea}


\section{Aceleración centrípeta}

La aceleración asociada al cambio de dirección del vector velocidad recibe el
nombre de aceleración centrípeta.

\begin{definicion}[Aceleración centrípeta]

La aceleración centrípeta es la componente de la aceleración que apunta hacia
el centro de curvatura de la trayectoria y describe el cambio de dirección de
la velocidad.

\end{definicion}

En una trayectoria circular:

\[
\boxed{
a_c=\frac{v^2}{R}
}.
\]

Como:

\[
v=R|\omega|,
\]

obtenemos:

\[
a_c
=
\frac{R^2\omega^2}{R}.
\]

Por tanto:

\[
\boxed{
a_c=R\omega^2
}.
\]

La aceleración centrípeta es una magnitud positiva:

\[
a_c\geq0.
\]

Su dirección apunta siempre hacia el centro.

\begin{advertencia}

El símbolo:

\[
a_c
\]

representa aquí la magnitud de la aceleración centrípeta.

No debe confundirse con el vector completo:

\[
\vect{a}_c.
\]

La magnitud no posee signo negativo, pero el vector sí tiene dirección y
sentido.

\end{advertencia}


\section{Representación vectorial de la aceleración centrípeta}

Consideremos una partícula en una trayectoria circular de radio \(R\) cuyo
centro coincide con el origen.

El vector de posición es:

\[
\vect{r}
=
x\ihat+y\jhat.
\]

La aceleración centrípeta apunta en sentido contrario al vector radial.

Por tanto:

\[
\boxed{
\vect{a}_c
=
-\frac{v^2}{R^2}\vect{r}
}.
\]

Como:

\[
v^2=R^2\omega^2,
\]

tenemos:

\[
\boxed{
\vect{a}_c=-\omega^2\vect{r}
}.
\]

En componentes rectangulares:

\[
\boxed{
a_{cx}=-\omega^2x
},
\]

\[
\boxed{
a_{cy}=-\omega^2y
}.
\]

Estas relaciones son válidas también para movimiento circular no uniforme
cuando describen únicamente la componente centrípeta.

\begin{interpretacion}

El signo negativo indica que la aceleración centrípeta es opuesta al vector
radial dirigido desde el centro hacia la partícula.

No significa que la aceleración sea siempre negativa en ambos ejes.

Los signos de sus componentes dependen de la posición de la partícula.

\end{interpretacion}


\section{Diagrama de velocidad y aceleración centrípeta}

\begin{center}
\begin{tikzpicture}[scale=1.0]

    \coordinate (O) at (0,0);

    \draw[-{Latex},thick]
        (-3.5,0) -- (3.6,0)
        node[right] {\(x\)};

    \draw[-{Latex},thick]
        (0,-3.5) -- (0,3.6)
        node[above] {\(y\)};

    \draw[very thick,AzulPrincipal]
        (O) circle (2.5);

    \fill (O) circle (2pt);
    \node[below left] at (O) {\(O\)};

    \foreach \ang in {0,90,180,270}{

        \pgfmathsetmacro{\px}{2.5*cos(\ang)}
        \pgfmathsetmacro{\py}{2.5*sin(\ang)}

        \fill (\px,\py) circle (2.5pt);

        \draw[-{Latex},very thick,VerdePrincipal]
            (\px,\py) --
            ++({-0.65*sin(\ang)},{0.65*cos(\ang)});

        \draw[-{Latex},very thick,RojoPrincipal]
            (\px,\py) --
            ++({-0.75*cos(\ang)},{-0.75*sin(\ang)});
    }

    \node[VerdePrincipal] at (3.2,0.8) {\(\vect{v}\)};
    \node[RojoPrincipal] at (1.3,0.35) {\(\vect{a}_c\)};

    \node[VerdePrincipal] at (-1.0,3.1) {\(\vect{v}\)};
    \node[RojoPrincipal] at (0.65,1.4) {\(\vect{a}_c\)};

\end{tikzpicture}
\end{center}

En cada punto:

\[
\vect{v}
\]

es tangente a la circunferencia, mientras:

\[
\vect{a}_c
\]

apunta hacia el centro.

Por tanto:

\[
\boxed{
\vect{v}\pdot\vect{a}_c=0
}.
\]

Es decir, ambos vectores son perpendiculares.


\section{Deducción geométrica de la aceleración centrípeta}

Para comprender de dónde surge:

\[
a_c=\frac{v^2}{R},
\]

consideremos una partícula que describe una circunferencia con rapidez
constante.

En dos instantes próximos, sus velocidades tienen magnitud \(v\), pero
diferentes direcciones.

Sean:

\[
\vect{v}_1
\]

y:

\[
\vect{v}_2.
\]

Durante el intervalo \(\Delta t\), el vector velocidad gira un ángulo pequeño
\(\Delta\theta\).

Como los vectores velocidad tienen la misma magnitud, el triángulo que forman
sus extremos es semejante, en el límite de ángulos pequeños, al triángulo
formado por los radios correspondientes.

Por tanto:

\[
\frac{|\Delta\vect{v}|}{v}
\approx
\frac{\Delta s}{R}.
\]

Como la rapidez es constante:

\[
\Delta s=v\Delta t.
\]

Entonces:

\[
\frac{|\Delta\vect{v}|}{v}
\approx
\frac{v\Delta t}{R}.
\]

Despejando:

\[
|\Delta\vect{v}|
\approx
\frac{v^2\Delta t}{R}.
\]

Dividimos entre \(\Delta t\):

\[
\frac{|\Delta\vect{v}|}{\Delta t}
\approx
\frac{v^2}{R}.
\]

En el límite:

\[
\Delta t\to0,
\]

la aproximación geométrica se convierte en una igualdad.

Por tanto:

\[
\boxed{
a_c=\frac{v^2}{R}
}.
\]

\begin{interpretacion}

La aceleración centrípeta aumenta cuando aumenta la rapidez.

En particular:

\[
a_c\propto v^2
\]

si el radio permanece constante.

Por ello, duplicar la rapidez cuadruplica la aceleración centrípeta.

También:

\[
a_c\propto\frac{1}{R}
\]

cuando la rapidez se mantiene fija.

Una curva más cerrada exige un cambio más rápido de dirección del vector
velocidad.

\end{interpretacion}


\section{Deducción mediante componentes rectangulares}

Podemos obtener el mismo resultado usando las ecuaciones cartesianas.

Para movimiento circular uniforme:

\[
\theta(t)=\theta_0+\omega t.
\]

La posición es:

\[
x(t)=R\cos(\theta_0+\omega t),
\]

\[
y(t)=R\sen(\theta_0+\omega t).
\]

Derivando:

\[
v_x=-R\omega\sen(\theta_0+\omega t),
\]

\[
v_y=R\omega\cos(\theta_0+\omega t).
\]

Derivamos nuevamente.

Para \(x\):

\[
a_x
=
-R\omega^2\cos(\theta_0+\omega t).
\]

Para \(y\):

\[
a_y
=
-R\omega^2\sen(\theta_0+\omega t).
\]

Por tanto:

\[
\vect{a}
=
-R\omega^2
\left[
\cos\theta\,\ihat+\sen\theta\,\jhat
\right].
\]

Pero:

\[
\vect{r}
=
R
\left[
\cos\theta\,\ihat+\sen\theta\,\jhat
\right].
\]

Entonces:

\[
\boxed{
\vect{a}=-\omega^2\vect{r}
}.
\]

En movimiento circular uniforme, esta es toda la aceleración.

Su magnitud es:

\[
a
=
R\omega^2.
\]

Como:

\[
v=R|\omega|,
\]

resulta:

\[
\boxed{
a_c=\frac{v^2}{R}=R\omega^2
}.
\]

Esta deducción muestra que el resultado geométrico y el cálculo por
componentes rectangulares son completamente coherentes.


\section{Ecuaciones del movimiento circular uniforme}

Si:

\[
\omega=\text{constante},
\]

entonces:

\[
\deriv{\theta}{t}=\omega.
\]

Integrando:

\[
\boxed{
\theta(t)=\theta_0+\omega t
}.
\]

Las componentes de posición son:

\[
\boxed{
x(t)=R\cos(\theta_0+\omega t)
},
\]

\[
\boxed{
y(t)=R\sen(\theta_0+\omega t)
}.
\]

Las componentes de velocidad:

\[
\boxed{
v_x(t)=-R\omega\sen(\theta_0+\omega t)
},
\]

\[
\boxed{
v_y(t)=R\omega\cos(\theta_0+\omega t)
}.
\]

Las componentes de aceleración:

\[
\boxed{
a_x(t)=-R\omega^2\cos(\theta_0+\omega t)
},
\]

\[
\boxed{
a_y(t)=-R\omega^2\sen(\theta_0+\omega t)
}.
\]

\begin{condicion}

Estas ecuaciones describen MCU alrededor de un centro fijo situado en el
origen, con radio y velocidad angular algebraica constantes.

Si el centro está en:

\[
(x_C,y_C),
\]

las ecuaciones de posición se modifican a:

\[
x=x_C+R\cos\theta,
\]

\[
y=y_C+R\sen\theta.
\]

\end{condicion}


\section{¿Qué magnitudes permanecen constantes en el MCU?}

En el MCU permanecen constantes:

\[
R,
\]

\[
v,
\]

\[
|\omega|,
\]

\[
a_c.
\]

También el período y la frecuencia:

\[
T=\text{constante},
\]

\[
f=\text{constante}.
\]

Sin embargo, los vectores:

\[
\vect{r},
\qquad
\vect{v},
\qquad
\vect{a}
\]

cambian continuamente de dirección.

\begin{errorfrecuente}

No debe afirmarse que la aceleración vectorial es constante en MCU.

La magnitud:

\[
a_c=\frac{v^2}{R}
\]

permanece constante, pero su dirección cambia continuamente para seguir
apuntando hacia el centro de la trayectoria.

Por tanto:

\[
\vect{a}\neq\text{constante}.
\]

\end{errorfrecuente}


\section{Gráfica de posición angular contra tiempo}

Para MCU:

\[
\theta(t)=\theta_0+\omega t.
\]

Esta ecuación corresponde a una recta.

En una gráfica de posición angular acumulada contra tiempo:

\[
\boxed{
\text{pendiente}=\omega
}.
\]

\begin{center}
\begin{tikzpicture}
\begin{axis}[
    width=0.80\textwidth,
    height=6cm,
    axis lines=left,
    xlabel={\(t\) (\si{\second})},
    ylabel={\(\theta\) (\si{\radian})},
    xmin=0,
    xmax=5,
    ymin=0,
    ymax=12,
    xtick={0,1,2,3,4,5},
    ytick={0,2,4,6,8,10,12},
    grid=both
]

\addplot[
    AzulPrincipal,
    very thick,
    domain=0:5,
    samples=2
]
{1+2*x};

\end{axis}
\end{tikzpicture}
\end{center}

En esta gráfica:

\[
\theta_0=\SI{1}{\radian},
\]

y:

\[
\omega=\SI{2}{\radian\per\second}.
\]

La pendiente positiva indica giro antihorario según nuestra convención.


\section{Gráficas de posición cartesiana en el MCU}

Aunque:

\[
\theta(t)
\]

cambia linealmente con el tiempo, las coordenadas cartesianas no lo hacen.

En MCU:

\[
x(t)=R\cos(\theta_0+\omega t),
\]

\[
y(t)=R\sen(\theta_0+\omega t).
\]

Por tanto, las gráficas:

\[
x-t
\]

y:

\[
y-t
\]

son funciones sinusoidales.

\begin{interpretacion}

Una trayectoria circular puede producir gráficas de posición cartesiana
oscilatorias.

Esto no significa que la partícula se mueva alternativamente sobre una línea
recta.

Significa que cada proyección cartesiana de su posición cambia de manera
periódica mientras la partícula recorre la circunferencia.

\end{interpretacion}


\section{Ejemplo: movimiento circular uniforme}

\begin{ejemplo}

Una partícula describe una circunferencia de radio:

\[
R=\SI{2.00}{\meter}
\]

con frecuencia:

\[
f=\SI{0.500}{\hertz}.
\]

En:

\[
t=0,
\]

se encuentra sobre el eje \(+x\) y gira en sentido antihorario.

Determine:

\begin{enumerate}[label=\alph*)]
    \item el período;
    \item la velocidad angular;
    \item la rapidez;
    \item la aceleración centrípeta;
    \item los vectores posición, velocidad y aceleración en:
    \[
    t=\SI{0.500}{\second}.
    \]
\end{enumerate}

\end{ejemplo}

El período es:

\[
T=\frac{1}{f}.
\]

Entonces:

\[
\boxed{
T=\SI{2.00}{\second}
}.
\]

Como el giro es antihorario:

\[
\omega=2\pi f.
\]

Entonces:

\[
\omega
=
2\pi(\SI{0.500}{\hertz}).
\]

Por tanto:

\[
\boxed{
\omega=\pi\ \si{\radian\per\second}
}.
\]

La rapidez es:

\[
v=R\omega.
\]

Así:

\[
v
=
(\SI{2.00}{\meter})
(\pi\ \si{\radian\per\second}).
\]

Por tanto:

\[
\boxed{
v\approx
\SI{6.28}{\meter\per\second}
}.
\]

La aceleración centrípeta es:

\[
a_c=R\omega^2.
\]

Entonces:

\[
a_c
=
2.00\pi^2
\ \si{\meter\per\second\squared}.
\]

Así:

\[
\boxed{
a_c\approx
\SI{19.7}{\meter\per\second\squared}
}.
\]

Ahora determinamos la posición angular.

Como:

\[
\theta_0=0,
\]

tenemos:

\[
\theta(t)=\omega t.
\]

Para:

\[
t=\SI{0.500}{\second},
\]

resulta:

\[
\theta
=
\pi(0.500)
=
\frac{\pi}{2}\ \si{\radian}.
\]

Entonces:

\[
\theta=90\grado.
\]

La partícula se encuentra en el punto superior de la circunferencia:

\[
x=R\cos90\grado=0,
\]

\[
y=R\sen90\grado=\SI{2.00}{\meter}.
\]

Por tanto:

\[
\boxed{
\vect{r}
=
2.00\jhat\ \si{\meter}
}.
\]

Las componentes de velocidad son:

\[
v_x=-R\omega\sen\theta,
\]

\[
v_y=R\omega\cos\theta.
\]

Entonces:

\[
v_x
=
-2.00\pi,
\]

\[
v_y=0.
\]

Así:

\[
\boxed{
\vect{v}
=
-6.28\ihat\,
\si{\meter\per\second}
}.
\]

La velocidad apunta hacia la izquierda, tangente a la trayectoria.

Finalmente:

\[
\vect{a}
=
-\omega^2\vect{r}.
\]

Entonces:

\[
\boxed{
\vect{a}
=
-19.7\jhat\,
\si{\meter\per\second\squared}
}.
\]

La aceleración apunta verticalmente hacia abajo, es decir, hacia el centro.


\section{Movimiento circular no uniforme}

Hasta ahora hemos estudiado el caso en el que la rapidez permanece constante.

Pero un objeto también puede moverse sobre una circunferencia mientras
aumenta o disminuye su rapidez.

Por ejemplo, al poner en marcha un ventilador, los puntos de sus aspas
comienzan a girar cada vez más rápido.

Cuando se apaga y disminuye su rapidez angular, ocurre lo contrario.

\begin{definicion}[Movimiento circular no uniforme]

El movimiento circular no uniforme es aquel en el que una partícula describe
una trayectoria circular, pero su rapidez no permanece constante.

\end{definicion}

En este caso el vector velocidad puede cambiar simultáneamente:

\begin{itemize}
    \item de magnitud;
    \item de dirección.
\end{itemize}

Por tanto, necesitamos distinguir dos contribuciones a la aceleración.


\section{Aceleración angular media}

La velocidad angular puede cambiar con el tiempo.

Para medir ese cambio introducimos la aceleración angular.

\begin{definicion}[Aceleración angular media]

La aceleración angular media es el cambio de velocidad angular dividido entre
el intervalo de tiempo:

\[
\boxed{
\bar{\alpha}
=
\frac{\Delta\omega}{\Delta t}
}.
\]

\end{definicion}

La unidad SI es:

\[
\si{\radian\per\second\squared}.
\]

Si:

\[
\bar{\alpha}>0,
\]

la velocidad angular algebraica aumenta.

Si:

\[
\bar{\alpha}<0,
\]

disminuye.

\begin{advertencia}

Una aceleración angular negativa no significa necesariamente que el objeto
esté disminuyendo su rapidez de giro.

La interpretación depende del signo de \(\omega\).

Si:

\[
\omega<0,
\qquad
\alpha<0,
\]

la magnitud de la velocidad angular aumenta.

\end{advertencia}


\section{Aceleración angular instantánea}

\begin{definicion}[Aceleración angular instantánea]

La aceleración angular instantánea es:

\[
\boxed{
\alpha
=
\deriv{\omega}{t}
}.
\]

Como:

\[
\omega=\deriv{\theta}{t},
\]

también:

\[
\boxed{
\alpha=\derivdos{\theta}{t}
}.
\]

\end{definicion}

En el sistema elegido:

\[
\vect{\alpha}
=
\alpha\khat.
\]

La aceleración angular indica qué tan rápidamente cambia la velocidad angular.

No debe confundirse con la aceleración lineal de la partícula.


\section{Aceleración tangencial}

Supongamos que la rapidez de una partícula cambia mientras recorre una
circunferencia.

La componente de aceleración responsable de ese cambio recibe el nombre de
aceleración tangencial.

\begin{definicion}[Aceleración tangencial]

La aceleración tangencial es la componente de la aceleración paralela a la
tangente de la trayectoria y responsable del cambio de rapidez.

\end{definicion}

Definimos su componente algebraica respecto de la tangente positiva como:

\[
\boxed{
a_t=\deriv{v_{\mathrm{tan}}}{t}
}.
\]

Como:

\[
v_{\mathrm{tan}}=R\omega,
\]

y el radio es constante:

\[
a_t
=
R\deriv{\omega}{t}.
\]

Por tanto:

\[
\boxed{
a_t=R\alpha
}.
\]

Su magnitud es:

\[
\boxed{
|a_t|=R|\alpha|
}.
\]

\begin{interpretacion}

La aceleración tangencial modifica la magnitud de la velocidad.

Si apunta en el mismo sentido que la velocidad, la rapidez aumenta.

Si apunta en sentido contrario, la rapidez disminuye.

\end{interpretacion}


\section{Aceleración centrípeta en movimiento circular no uniforme}

Aunque la rapidez cambie, la partícula continúa moviéndose sobre una
circunferencia.

Por ello, la dirección del vector velocidad sigue cambiando.

La componente centrípeta continúa existiendo:

\[
\boxed{
a_c=\frac{v^2}{R}
}.
\]

Equivalentemente:

\[
\boxed{
a_c=R\omega^2
}.
\]

El movimiento circular no uniforme presenta entonces dos componentes de
aceleración:

\begin{itemize}
    \item aceleración tangencial;
    \item aceleración centrípeta.
\end{itemize}

No son dos aceleraciones de partículas diferentes.

Son componentes de un mismo vector aceleración.


\section{Aceleración total}

En una trayectoria circular, la aceleración total puede expresarse como:

\[
\boxed{
\vect{a}
=
\vect{a}_t+\vect{a}_c
}.
\]

La aceleración tangencial es paralela a la trayectoria.

La aceleración centrípeta apunta hacia el centro.

Ambas componentes son perpendiculares.

Por tanto, la magnitud de la aceleración total es:

\[
\boxed{
a=\sqrt{a_t^2+a_c^2}
}.
\]

Sustituyendo sus magnitudes:

\[
\boxed{
a
=
\sqrt{
R^2\alpha^2+R^2\omega^4
}
}.
\]

Factorizando:

\[
\boxed{
a
=
R\sqrt{\alpha^2+\omega^4}
}.
\]

\begin{condicion}

Esta expresión requiere movimiento circular de radio constante alrededor de
un centro fijo.

Las magnitudes \(\alpha\) y \(\omega\) deben expresarse en unidades coherentes
del SI.

\end{condicion}


\section{Representación geométrica de las aceleraciones}

\begin{center}
\begin{tikzpicture}[scale=1.0]

    \coordinate (O) at (0,0);
    \coordinate (P) at (2.5,0);

    \draw[-{Latex},thick]
        (-3.3,0) -- (3.5,0)
        node[right] {\(x\)};

    \draw[-{Latex},thick]
        (0,-3.3) -- (0,3.4)
        node[above] {\(y\)};

    \draw[very thick,AzulPrincipal]
        (O) circle (2.5);

    \fill (O) circle (2pt);
    \node[below left] at (O) {\(O\)};

    \fill (P) circle (3pt);
    \node[below right] at (P) {\(P\)};

    \draw[-{Latex},very thick,VerdePrincipal]
        (P) -- ++(0,1.55)
        node[above] {\(\vect{v}\)};

    \draw[-{Latex},very thick,RojoPrincipal]
        (P) -- ++(-1.55,0)
        node[midway,below] {\(\vect{a}_c\)};

    \draw[-{Latex},very thick,AzulSecundario]
        (P) -- ++(0,0.85)
        node[right] {\(\vect{a}_t\)};

    \draw[-{Latex},thick,dashed]
        (P) -- ++(-1.55,0.85)
        node[above left] {\(\vect{a}\)};

\end{tikzpicture}
\end{center}

El dibujo representa un instante en el que la partícula se mueve hacia arriba
y aumenta su rapidez.

En ese instante:

\[
\vect{a}_t
\]

apunta hacia arriba, mientras:

\[
\vect{a}_c
\]

apunta hacia el centro.

La aceleración total tiene una dirección intermedia.

Si la partícula estuviera disminuyendo su rapidez, la componente tangencial
apuntaría en sentido contrario al vector velocidad.


\section{Deducción de la aceleración total por componentes rectangulares}

La posición sigue siendo:

\[
\vect{r}
=
R\cos\theta\,\ihat
+
R\sen\theta\,\jhat.
\]

Pero ahora:

\[
\omega=\omega(t).
\]

La velocidad es:

\[
\vect{v}
=
-R\omega\sen\theta\,\ihat
+
R\omega\cos\theta\,\jhat.
\]

Derivamos cada componente.

Para \(x\):

\[
a_x
=
\deriv{}{t}
(-R\omega\sen\theta).
\]

Aplicando la regla del producto:

\[
a_x
=
-R\alpha\sen\theta
-
R\omega^2\cos\theta.
\]

Para \(y\):

\[
a_y
=
\deriv{}{t}
(R\omega\cos\theta).
\]

Entonces:

\[
a_y
=
R\alpha\cos\theta
-
R\omega^2\sen\theta.
\]

Por tanto:

\[
\boxed{
a_x
=
-R\alpha\sen\theta
-
R\omega^2\cos\theta
},
\]

\[
\boxed{
a_y
=
R\alpha\cos\theta
-
R\omega^2\sen\theta
}.
\]

El vector aceleración total es:

\[
\boxed{
\vect{a}
=
\left(
-R\alpha\sen\theta-R\omega^2\cos\theta
\right)\ihat
+
\left(
R\alpha\cos\theta-R\omega^2\sen\theta
\right)\jhat
}.
\]

Podemos separar dos contribuciones:

\[
\vect{a}_t
=
R\alpha
\left(
-\sen\theta\,\ihat+\cos\theta\,\jhat
\right),
\]

\[
\vect{a}_c
=
-R\omega^2
\left(
\cos\theta\,\ihat+\sen\theta\,\jhat
\right).
\]

La primera es tangencial.

La segunda es radial hacia el centro.

\begin{interpretacion}

El análisis por componentes rectangulares demuestra que la aceleración total
de una partícula en movimiento circular puede cambiar tanto la magnitud como
la dirección de la velocidad.

No necesitamos abandonar las coordenadas cartesianas para describir el
movimiento circular.

Las magnitudes angulares permiten simplificar el análisis, pero las
componentes rectangulares siguen proporcionando una descripción general y
completa.

\end{interpretacion}


\section{Movimiento circular uniformemente acelerado}

Un caso particular del movimiento circular no uniforme ocurre cuando la
aceleración angular permanece constante.

\begin{definicion}[Movimiento circular con aceleración angular constante]

Es el movimiento circular en el que:

\[
\boxed{
\alpha=\text{constante}
}.
\]

\end{definicion}

También se conoce como movimiento circular uniformemente variado.

Si:

\[
\alpha>0,
\]

la velocidad angular algebraica aumenta.

Si:

\[
\alpha<0,
\]

disminuye.

Esto no implica necesariamente que aumente o disminuya la rapidez angular,
pues debemos considerar también el signo de \(\omega\).


\section{Ecuación de velocidad angular}

Partimos de:

\[
\alpha=\deriv{\omega}{t}.
\]

Si:

\[
\alpha=\text{constante},
\]

entonces:

\[
\dd\omega=\alpha\,\dd t.
\]

Integramos desde:

\[
t=0,
\qquad
\omega=\omega_0,
\]

hasta un tiempo \(t\):

\[
\int_{\omega_0}^{\omega}\dd\omega
=
\int_0^t\alpha\,\dd t.
\]

Por tanto:

\[
\omega-\omega_0=\alpha t.
\]

Así:

\[
\boxed{
\omega=\omega_0+\alpha t
}.
\]

\begin{interpretacion}

La velocidad angular cambia linealmente con el tiempo.

La cantidad:

\[
\alpha t
\]

representa el cambio angular de velocidad acumulado durante el intervalo.

\end{interpretacion}


\section{Ecuación de posición angular}

Sabemos que:

\[
\omega=\deriv{\theta}{t}.
\]

Como:

\[
\omega=\omega_0+\alpha t,
\]

tenemos:

\[
\deriv{\theta}{t}
=
\omega_0+\alpha t.
\]

Integramos:

\[
\theta-\theta_0
=
\int_0^t
(\omega_0+\alpha t)\,\dd t.
\]

Entonces:

\[
\boxed{
\theta
=
\theta_0
+
\omega_0t
+
\frac{1}{2}\alpha t^2
}.
\]

Equivalentemente:

\[
\boxed{
\Delta\theta
=
\omega_0t
+
\frac{1}{2}\alpha t^2
}.
\]

Esta ecuación es análoga a la ecuación de posición del MRUA:

\[
x=x_0+v_0t+\frac{1}{2}at^2.
\]

Sin embargo, las magnitudes poseen significados físicos diferentes.


\section{Relación angular sin el tiempo}

Partimos de:

\[
\omega=\omega_0+\alpha t.
\]

Para aceleración angular constante:

\[
\Delta\theta
=
\frac{\omega_0+\omega}{2}t.
\]

Despejando:

\[
t=\frac{\omega-\omega_0}{\alpha},
\]

para:

\[
\alpha\neq0.
\]

Sustituimos:

\[
\Delta\theta
=
\frac{\omega_0+\omega}{2}
\frac{\omega-\omega_0}{\alpha}.
\]

Entonces:

\[
2\alpha\Delta\theta
=
(\omega_0+\omega)(\omega-\omega_0).
\]

Por diferencia de cuadrados:

\[
2\alpha\Delta\theta
=
\omega^2-\omega_0^2.
\]

Por tanto:

\[
\boxed{
\omega^2
=
\omega_0^2
+
2\alpha\Delta\theta
}.
\]

\begin{condicion}

Esta ecuación requiere aceleración angular constante.

Como contiene velocidades angulares al cuadrado, no determina por sí sola el
sentido de giro.

Después de resolver, debe recuperarse el signo físico de \(\omega\).

\end{condicion}


\section{Ecuaciones angulares con aceleración constante}

Las relaciones principales son:

\[
\boxed{
\omega=\omega_0+\alpha t
},
\]

\[
\boxed{
\theta
=
\theta_0+\omega_0t+\frac{1}{2}\alpha t^2
},
\]

\[
\boxed{
\omega^2
=
\omega_0^2+2\alpha\Delta\theta
},
\]

\[
\boxed{
\Delta\theta
=
\frac{\omega_0+\omega}{2}t
}.
\]

Estas ecuaciones son válidas únicamente cuando:

\[
\alpha=\text{constante}.
\]


\section{Comparación entre movimientos lineales y angulares}

Existe una analogía entre las ecuaciones del movimiento rectilíneo con
aceleración constante y las ecuaciones angulares con aceleración constante.

\begin{center}
\begin{tabularx}{\textwidth}{>{\raggedright\arraybackslash}X
                            >{\raggedright\arraybackslash}X}
\toprule
Movimiento lineal & Movimiento angular \\
\midrule
Posición \(x\) & Posición angular \(\theta\) \\

Velocidad \(v_x\) & Velocidad angular \(\omega\) \\

Aceleración \(a_x\) & Aceleración angular \(\alpha\) \\

\(\Delta x=v_{0x}t+\frac12a_xt^2\)
&
\(\Delta\theta=\omega_0t+\frac12\alpha t^2\)
\\

\(v_x=v_{0x}+a_xt\)
&
\(\omega=\omega_0+\alpha t\)
\\

\(v_x^2=v_{0x}^2+2a_x\Delta x\)
&
\(\omega^2=\omega_0^2+2\alpha\Delta\theta\)
\\
\bottomrule
\end{tabularx}
\end{center}

\begin{advertencia}

Esta analogía no significa que la aceleración lineal total de una partícula en
movimiento circular sea constante cuando \(\alpha\) es constante.

Incluso con:

\[
\alpha=\text{constante},
\]

la aceleración centrípeta:

\[
a_c=R\omega^2
\]

puede cambiar de magnitud porque cambia \(\omega\).

Además, las direcciones de las componentes de aceleración cambian mientras
la partícula gira.

\end{advertencia}


\section{Ejemplo: rueda que aumenta su rapidez angular}

\begin{ejemplo}

Una rueda parte del reposo y gira con aceleración angular constante:

\[
\alpha=\SI{2.00}{\radian\per\second\squared}.
\]

Determine después de:

\[
t=\SI{4.00}{\second}
\]

\begin{enumerate}[label=\alph*)]
    \item la velocidad angular;
    \item el desplazamiento angular;
    \item el número de revoluciones completadas.
\end{enumerate}

\end{ejemplo}

Como parte del reposo:

\[
\omega_0=0.
\]

La velocidad angular es:

\[
\omega=\omega_0+\alpha t.
\]

Sustituyendo:

\[
\omega
=
0+
(\SI{2.00}{\radian\per\second\squared})
(\SI{4.00}{\second}).
\]

Entonces:

\[
\boxed{
\omega=\SI{8.00}{\radian\per\second}
}.
\]

El desplazamiento angular es:

\[
\Delta\theta
=
\omega_0t+\frac12\alpha t^2.
\]

Como:

\[
\omega_0=0,
\]

resulta:

\[
\Delta\theta
=
\frac12(2.00)(4.00)^2.
\]

Por tanto:

\[
\boxed{
\Delta\theta=\SI{16.0}{\radian}
}.
\]

Una revolución equivale a:

\[
2\pi\ \si{\radian}.
\]

Entonces:

\[
N
=
\frac{\Delta\theta}{2\pi}.
\]

Así:

\[
N
=
\frac{16.0}{2\pi}
\approx2.55.
\]

Por tanto:

\[
\boxed{N\approx2.55\text{ revoluciones}}.
\]

La rueda ha completado dos vueltas enteras y una fracción adicional de vuelta.


\section{Ejemplo: aceleración tangencial y centrípeta}

\begin{ejemplo}

Una partícula se mueve sobre una circunferencia de radio:

\[
R=\SI{0.80}{\meter}.
\]

En cierto instante tiene:

\[
\omega=\SI{5.0}{\radian\per\second}
\]

y:

\[
\alpha=\SI{2.0}{\radian\per\second\squared}.
\]

Determine:

\begin{enumerate}[label=\alph*)]
    \item la rapidez;
    \item la magnitud de la aceleración tangencial;
    \item la magnitud de la aceleración centrípeta;
    \item la magnitud de la aceleración total.
\end{enumerate}

\end{ejemplo}

La rapidez es:

\[
v=R|\omega|.
\]

Entonces:

\[
v
=
(0.80)(5.0).
\]

Así:

\[
\boxed{
v=\SI{4.0}{\meter\per\second}
}.
\]

La aceleración tangencial tiene magnitud:

\[
|a_t|=R|\alpha|.
\]

Por tanto:

\[
|a_t|
=
(0.80)(2.0).
\]

Entonces:

\[
\boxed{
|a_t|=\SI{1.6}{\meter\per\second\squared}
}.
\]

La aceleración centrípeta es:

\[
a_c=R\omega^2.
\]

Entonces:

\[
a_c
=
(0.80)(5.0)^2.
\]

Por tanto:

\[
\boxed{
a_c=\SI{20}{\meter\per\second\squared}
}.
\]

Como las componentes tangencial y centrípeta son perpendiculares:

\[
a
=
\sqrt{a_t^2+a_c^2}.
\]

Entonces:

\[
a
=
\sqrt{(1.6)^2+(20)^2}.
\]

Así:

\[
\boxed{
a\approx\SI{20.1}{\meter\per\second\squared}
}.
\]

\begin{interpretacion}

En este instante la aceleración centrípeta es mucho mayor que la tangencial.

Por ello, la aceleración total apunta principalmente hacia el centro, aunque
presenta también una componente tangencial asociada al aumento de rapidez.

\end{interpretacion}


\section{Velocidad y aceleración en un punto determinado}

En problemas más completos no basta conocer las magnitudes.

También necesitamos obtener las componentes rectangulares.

Supongamos que una partícula se encuentra en una posición angular \(\theta\).

Entonces:

\[
x=R\cos\theta,
\qquad
y=R\sen\theta.
\]

Las componentes de velocidad son:

\[
v_x=-R\omega\sen\theta,
\]

\[
v_y=R\omega\cos\theta.
\]

Las componentes de aceleración total son:

\[
a_x
=
-R\alpha\sen\theta-R\omega^2\cos\theta,
\]

\[
a_y
=
R\alpha\cos\theta-R\omega^2\sen\theta.
\]

Estas ecuaciones permiten conocer el vector velocidad y el vector aceleración
en cualquier punto de la trayectoria, siempre que se conozcan:

\[
R,\qquad
\theta,\qquad
\omega,\qquad
\alpha.
\]


\section{Ejemplo: análisis mediante componentes rectangulares}

\begin{ejemplo}

Una partícula describe una circunferencia de radio:

\[
R=\SI{2.0}{\meter}.
\]

En cierto instante su posición angular es:

\[
\theta=60\grado,
\]

su velocidad angular es:

\[
\omega=\SI{3.0}{\radian\per\second},
\]

y su aceleración angular es:

\[
\alpha=-\SI{1.0}{\radian\per\second\squared}.
\]

Determine los vectores posición, velocidad y aceleración en ese instante.

\end{ejemplo}

La posición es:

\[
\vect{r}
=
R\cos\theta\,\ihat
+
R\sen\theta\,\jhat.
\]

Sustituyendo:

\[
x=2.0\cos60\grado,
\]

\[
y=2.0\sen60\grado.
\]

Por tanto:

\[
x=\SI{1.00}{\meter},
\]

\[
y\approx\SI{1.73}{\meter}.
\]

Así:

\[
\boxed{
\vect{r}
=
(1.00\ihat+1.73\jhat)\,\si{\meter}
}.
\]

Para la velocidad:

\[
v_x=-R\omega\sen\theta.
\]

Entonces:

\[
v_x
=
-(2.0)(3.0)\sen60\grado.
\]

Por tanto:

\[
v_x\approx
-\SI{5.20}{\meter\per\second}.
\]

La componente vertical es:

\[
v_y=R\omega\cos\theta.
\]

Entonces:

\[
v_y
=
(2.0)(3.0)\cos60\grado.
\]

Así:

\[
v_y=\SI{3.00}{\meter\per\second}.
\]

Por tanto:

\[
\boxed{
\vect{v}
=
(-5.20\ihat+3.00\jhat)
\,\si{\meter\per\second}
}.
\]

Ahora analizamos la aceleración.

La componente horizontal es:

\[
a_x
=
-R\alpha\sen\theta-R\omega^2\cos\theta.
\]

Sustituyendo:

\[
a_x
=
-(2.0)(-1.0)\sen60\grado
-
(2.0)(3.0)^2\cos60\grado.
\]

Entonces:

\[
a_x
\approx
1.73-9.00.
\]

Por tanto:

\[
a_x
\approx
-\SI{7.27}{\meter\per\second\squared}.
\]

La componente vertical es:

\[
a_y
=
R\alpha\cos\theta-R\omega^2\sen\theta.
\]

Entonces:

\[
a_y
=
(2.0)(-1.0)\cos60\grado
-
(2.0)(3.0)^2\sen60\grado.
\]

Así:

\[
a_y
\approx
-1.00-15.59.
\]

Por tanto:

\[
a_y
\approx
-\SI{16.6}{\meter\per\second\squared}.
\]

Finalmente:

\[
\boxed{
\vect{a}
=
(-7.27\ihat-16.6\jhat)
\,\si{\meter\per\second\squared}
}.
\]

\begin{interpretacion}

El vector aceleración posee componentes negativas en ambos ejes.

Sin embargo, no debemos interpretar esto como una simple aceleración
``hacia abajo y hacia la izquierda'' sin considerar la geometría del
movimiento.

El resultado combina:

\begin{itemize}
    \item una componente centrípeta hacia el centro;
    \item una componente tangencial opuesta al sentido del movimiento.
\end{itemize}

Como:

\[
\omega>0,
\qquad
\alpha<0,
\]

la partícula gira en sentido antihorario, pero su rapidez está disminuyendo.

\end{interpretacion}


\section{Gráficas del movimiento circular con aceleración angular constante}

Cuando:

\[
\alpha=\text{constante},
\]

la velocidad angular cambia linealmente:

\[
\omega=\omega_0+\alpha t.
\]

Por ello, la gráfica:

\[
\omega-t
\]

es una recta.

Su pendiente es:

\[
\boxed{
\text{pendiente}=\alpha
}.
\]

El área algebraica bajo la gráfica \(\omega-t\) representa:

\[
\boxed{
\Delta\theta
=
\int_{t_i}^{t_f}\omega(t)\,\dd t
}.
\]

La posición angular:

\[
\theta
=
\theta_0+\omega_0t+\frac12\alpha t^2
\]

es una función cuadrática del tiempo.

Por tanto, su gráfica:

\[
\theta-t
\]

es parabólica cuando:

\[
\alpha\neq0.
\]

Finalmente, la gráfica:

\[
\alpha-t
\]

es una recta horizontal cuando la aceleración angular es constante.


\section{Cambio de sentido en el movimiento circular}

Una partícula puede detener momentáneamente su giro y comenzar a moverse en
sentido contrario.

Esto ocurre cuando la velocidad angular pasa por cero:

\[
\omega=0.
\]

Si la aceleración angular es constante:

\[
\omega=\omega_0+\alpha t.
\]

Entonces:

\[
0=\omega_0+\alpha t.
\]

Por tanto:

\[
\boxed{
t_{\mathrm{cambio}}
=
-\frac{\omega_0}{\alpha}
}.
\]

Esta expresión proporciona un instante posterior al inicial solamente cuando:

\[
t_{\mathrm{cambio}}>0.
\]

\begin{interpretacion}

En el instante en que:

\[
\omega=0,
\]

la rapidez lineal también es cero:

\[
v=R|\omega|=0.
\]

La aceleración centrípeta desaparece momentáneamente:

\[
a_c=R\omega^2=0.
\]

Sin embargo, la aceleración tangencial puede continuar siendo distinta de
cero si:

\[
\alpha\neq0.
\]

Por tanto, el reposo instantáneo no implica aceleración nula.

\end{interpretacion}


\section{Movimiento circular y observadores}

La descripción de un movimiento circular depende del sistema de referencia.

Por ejemplo, un punto marcado en el borde de una rueda que gira alrededor
de un eje fijo describe una circunferencia respecto de un observador fijo al
eje.

Si esa rueda además se traslada sobre el suelo, el mismo punto puede
describir una trayectoria más compleja respecto de un observador situado en
la carretera.

Por tanto, es importante especificar si estamos estudiando:

\begin{itemize}
    \item el movimiento alrededor de un centro fijo;
    \item el movimiento respecto de un sistema que se traslada;
    \item la combinación de traslación y rotación.
\end{itemize}

Las ecuaciones desarrolladas en este capítulo corresponden principalmente al
primer caso.


\section{Aplicaciones del movimiento circular}

\begin{aplicacion}[Centrifugadoras]

Una centrifugadora hace girar muestras alrededor de un eje.

Los puntos situados a mayor distancia del eje y con la misma rapidez angular
poseen mayor rapidez lineal:

\[
v=R|\omega|.
\]

También presentan mayor aceleración centrípeta:

\[
a_c=R\omega^2.
\]

Estas relaciones permiten describir cinemáticamente el movimiento de las
muestras durante el funcionamiento del equipo.

\end{aplicacion}


\begin{aplicacion}[Ruedas y mecanismos giratorios]

En una rueda rígida que gira alrededor de un eje fijo, todos los puntos
comparten la misma velocidad angular.

Sin embargo, la rapidez lineal depende del radio.

Por ello, los puntos del borde se desplazan más rápidamente que los puntos
cercanos al centro.

\end{aplicacion}


\begin{aplicacion}[Movimiento de satélites]

En ciertos problemas resulta útil aproximar la trayectoria de un satélite como
circular.

Si el radio orbital y el período son conocidos, pueden estimarse la rapidez y
la aceleración centrípeta mediante:

\[
v=\frac{2\pi R}{T},
\]

\[
a_c=\frac{v^2}{R}.
\]

Esta descripción es cinemática: permite calcular cómo se mueve el satélite sin
determinar todavía la interacción responsable de mantener su trayectoria.

\end{aplicacion}


\section{Actividad experimental}

Una práctica sencilla consiste en estudiar un objeto que gira alrededor de un
eje fijo.

Puede utilizarse un disco con puntos marcados a diferentes distancias del
centro.

Se mide:

\begin{itemize}
    \item el radio de cada punto;
    \item el tiempo necesario para completar varias vueltas;
    \item el número total de vueltas;
    \item el sentido del movimiento.
\end{itemize}

Con los datos se determina:

\[
T=\frac{\Delta t}{N},
\]

\[
f=\frac{1}{T},
\]

\[
|\omega|=2\pi f.
\]

Para cada radio puede calcularse:

\[
v=R|\omega|.
\]

Si la velocidad angular es común a todos los puntos, una gráfica de rapidez
lineal contra radio debe ser una recta que pasa por el origen:

\[
v\propto R.
\]

Su pendiente corresponde a:

\[
|\omega|.
\]

De forma semejante:

\[
a_c=R\omega^2.
\]

Por tanto, para velocidad angular constante, una gráfica de \(a_c\) contra \(R\)
también es lineal, con pendiente:

\[
\omega^2.
\]

Esta actividad conecta las mediciones con las relaciones físicas desarrolladas.


\section{Estrategia general para resolver problemas}

\begin{estrategia}[Movimiento circular]

Para resolver un problema de movimiento circular:

\begin{enumerate}

    \item Identifique el objeto o punto cuya trayectoria se estudiará.

    \item Determine el centro y el radio de la circunferencia.

    \item Defina el sistema de referencia y el sentido angular positivo.

    \item Identifique si el movimiento es circular uniforme o no uniforme.

    \item Escriba los datos conocidos:
    \[
    R,\quad
    \theta_0,\quad
    \omega_0,\quad
    \alpha,\quad
    t.
    \]

    \item Convierta las magnitudes angulares a radianes cuando corresponda.

    \item Utilice las ecuaciones angulares apropiadas.

    \item Si se requieren velocidades o aceleraciones lineales, determine las
    componentes correspondientes.

    \item Para MCU:
    \[
    v=R|\omega|,
    \]
    \[
    a_c=R\omega^2.
    \]

    \item Para movimiento circular no uniforme:
    \[
    a_t=R\alpha,
    \]
    \[
    a_c=R\omega^2.
    \]

    \item Si se pide el vector aceleración, reconstruya sus componentes
    rectangulares.

    \item Verifique:
    \begin{itemize}
        \item unidades;
        \item signos;
        \item sentido de giro;
        \item dirección de la velocidad;
        \item dirección de la aceleración;
        \item condiciones de validez.
    \end{itemize}

\end{enumerate}

\end{estrategia}


\section{Errores frecuentes}

\begin{errorfrecuente}

\textbf{Creer que rapidez constante significa aceleración cero.}

En MCU:

\[
v=\text{constante},
\]

pero:

\[
a_c=\frac{v^2}{R}\neq0
\]

si:

\[
v\neq0.
\]

\end{errorfrecuente}


\begin{errorfrecuente}

\textbf{Confundir velocidad angular con velocidad lineal.}

La velocidad angular mide cambio de posición angular por unidad de tiempo.

La velocidad lineal describe el cambio del vector de posición en el espacio.

Sus unidades son diferentes:

\[
[\omega]=\si{\radian\per\second},
\]

\[
[v]=\si{\meter\per\second}.
\]

\end{errorfrecuente}


\begin{errorfrecuente}

\textbf{Creer que todos los puntos de una rueda poseen la misma rapidez lineal.}

Si giran alrededor del mismo eje con la misma rapidez angular:

\[
v=R|\omega|.
\]

Los puntos más alejados del eje tienen mayor rapidez lineal.

\end{errorfrecuente}


\begin{errorfrecuente}

\textbf{Confundir aceleración centrípeta con aceleración tangencial.}

La aceleración centrípeta cambia la dirección de la velocidad.

La aceleración tangencial cambia su magnitud.

Son componentes perpendiculares de la aceleración total.

\end{errorfrecuente}


\begin{errorfrecuente}

\textbf{Creer que la aceleración centrípeta apunta hacia afuera.}

La aceleración centrípeta apunta siempre hacia el centro de la trayectoria.

\end{errorfrecuente}


\begin{errorfrecuente}

\textbf{Afirmar que una aceleración angular negativa siempre reduce la rapidez.}

La rapidez de giro disminuye cuando \(\omega\) y \(\alpha\) tienen signos
opuestos.

Si poseen el mismo signo, la rapidez de giro aumenta.

\end{errorfrecuente}


\begin{errorfrecuente}

\textbf{Utilizar grados directamente en:}

\[
s=R\theta.
\]

El ángulo debe expresarse en radianes.

\end{errorfrecuente}


\begin{errorfrecuente}

\textbf{Confundir una revolución con desplazamiento vectorial.}

Después de una revolución completa:

\[
|\Delta\theta|=2\pi\ \si{\radian},
\]

pero:

\[
\Delta\vect{r}=\vect{0}.
\]

\end{errorfrecuente}


\begin{errorfrecuente}

\textbf{Confundir aceleración angular constante con aceleración lineal
vectorial constante.}

Aunque:

\[
\alpha=\text{constante},
\]

la aceleración lineal puede cambiar tanto de magnitud como de dirección.

\end{errorfrecuente}


\section{Secuencia sugerida para impartir la clase}

\begin{notapedagogica}

Una secuencia posible para desarrollar este capítulo es:

\begin{enumerate}

    \item Mostrar una rueda o un ventilador y preguntar cómo describir la
    posición de un punto de su borde.

    \item Recordar que la velocidad es tangente a la trayectoria.

    \item Establecer el sistema de referencia y el sentido positivo de giro.

    \item Introducir la posición angular y el radián.

    \item Relacionar arco y ángulo:
    \[
    s=R\theta.
    \]

    \item Diferenciar distancia recorrida, desplazamiento angular y
    desplazamiento vectorial.

    \item Introducir velocidad angular media e instantánea.

    \item Deducir:
    \[
    v=R|\omega|.
    \]

    \item Explicar período y frecuencia mediante vueltas reales.

    \item Definir movimiento circular uniforme.

    \item Comparar dos vectores velocidad en puntos distintos de una
    circunferencia.

    \item Establecer que existe aceleración aunque la rapidez sea constante.

    \item Deducir e interpretar:
    \[
    a_c=\frac{v^2}{R}.
    \]

    \item Mostrar los vectores de velocidad y aceleración sobre una
    circunferencia.

    \item Expresar posición, velocidad y aceleración mediante componentes
    rectangulares.

    \item Resolver un ejemplo numérico completo de MCU.

    \item Introducir movimiento circular no uniforme.

    \item Definir aceleración angular.

    \item Diferenciar aceleración tangencial y centrípeta.

    \item Construir la aceleración total mediante sus componentes.

    \item Deducir las ecuaciones angulares con aceleración constante.

    \item Resolver ejemplos de aumento y disminución de rapidez angular.

    \item Interpretar las gráficas:
    \[
    \theta-t,\qquad
    \omega-t,\qquad
    \alpha-t.
    \]

    \item Finalizar con preguntas que obliguen a distinguir magnitudes
    angulares, lineales, escalares y vectoriales.

\end{enumerate}

\end{notapedagogica}


\section{Preguntas de comprensión}

\begin{enumerate}

    \item ¿Qué característica geométrica define el movimiento circular?

    \item ¿Puede una partícula moverse circularmente con rapidez variable?

    \item ¿Qué representa el radio de la trayectoria?

    \item ¿Qué diferencia existe entre posición angular y desplazamiento
    angular?

    \item ¿Qué es un radián?

    \item ¿Por qué resulta conveniente utilizar radianes?

    \item ¿Cuántos radianes corresponden a una revolución completa?

    \item ¿Qué representa físicamente:
    \[
    s=R\theta?
    \]

    \item ¿Qué diferencia existe entre la longitud de arco y la magnitud del
    desplazamiento vectorial?

    \item Después de una vuelta completa, ¿cuál es el desplazamiento vectorial?

    \item ¿Qué representa la velocidad angular?

    \item ¿Cuál es la diferencia entre velocidad angular y rapidez lineal?

    \item ¿Por qué un punto del borde de una rueda posee mayor rapidez lineal
    que uno próximo al centro?

    \item ¿Qué significa el período?

    \item ¿Qué significa la frecuencia?

    \item ¿Por qué período y frecuencia son inversamente proporcionales?

    \item ¿Qué significa movimiento circular uniforme?

    \item ¿Por qué existe aceleración centrípeta en MCU?

    \item ¿En qué dirección apunta la velocidad instantánea?

    \item ¿En qué dirección apunta la aceleración centrípeta?

    \item ¿Por qué la velocidad y la aceleración centrípeta son
    perpendiculares?

    \item Si se duplica la rapidez lineal manteniendo el radio, ¿qué ocurre
    con la aceleración centrípeta?

    \item Si se duplica el radio manteniendo la rapidez angular, ¿qué ocurre
    con la aceleración centrípeta?

    \item ¿Qué diferencia existe entre la aceleración angular y la
    aceleración tangencial?

    \item ¿Qué significa que la aceleración angular sea negativa?

    \item ¿Puede un objeto tener aceleración angular negativa mientras
    aumenta su rapidez de giro?

    \item ¿Cuándo desaparece la aceleración tangencial?

    \item ¿Cuándo desaparece la aceleración centrípeta?

    \item ¿Puede existir aceleración centrípeta sin aceleración tangencial?

    \item ¿Puede una partícula tener aceleración tangencial distinta de cero
    y rapidez instantánea igual a cero?

    \item ¿Por qué la aceleración lineal total no es constante
    necesariamente cuando la aceleración angular sí lo es?

    \item ¿Qué representa la pendiente de una gráfica \(\theta-t\)?

    \item ¿Qué representa la pendiente de una gráfica \(\omega-t\)?

    \item ¿Qué representa el área algebraica bajo una gráfica
    \(\omega-t\)?

    \item ¿Cómo pueden determinarse las componentes rectangulares de la
    aceleración de una partícula en movimiento circular?

\end{enumerate}


\begin{resumen}

El movimiento circular es aquel en el que una partícula describe una
circunferencia de radio constante.

La posición angular permite describir la ubicación de la partícula alrededor
del centro:

\[
\theta=\theta(t).
\]

La medida angular natural para relacionar ángulos y longitudes es el radián:

\[
\boxed{s=R\theta}.
\]

El desplazamiento angular es:

\[
\boxed{
\Delta\theta=\theta_f-\theta_i
}.
\]

La velocidad angular es:

\[
\boxed{
\omega=\deriv{\theta}{t}
}.
\]

La aceleración angular es:

\[
\boxed{
\alpha=\deriv{\omega}{t}
}.
\]

La rapidez lineal se relaciona con la velocidad angular mediante:

\[
\boxed{
v=R|\omega|
}.
\]

En movimiento circular uniforme:

\[
R=\text{constante},
\]

\[
v=\text{constante},
\]

\[
\omega=\text{constante}.
\]

Sin embargo, la dirección del vector velocidad cambia continuamente.

Por ello existe aceleración centrípeta:

\[
\boxed{
a_c=\frac{v^2}{R}=R\omega^2
}.
\]

Esta aceleración apunta hacia el centro de la trayectoria.

En componentes rectangulares:

\[
\boxed{
\vect{a}_c=-\omega^2\vect{r}
}.
\]

Para movimiento circular uniforme también:

\[
\boxed{
|\omega|=\frac{2\pi}{T}=2\pi f
}.
\]

En movimiento circular no uniforme puede existir aceleración tangencial:

\[
\boxed{
a_t=R\alpha
}.
\]

La aceleración total combina las componentes centrípeta y tangencial:

\[
\boxed{
\vect{a}=\vect{a}_c+\vect{a}_t
}.
\]

Como ambas son perpendiculares:

\[
\boxed{
a=\sqrt{a_c^2+a_t^2}
}.
\]

Cuando la aceleración angular es constante:

\[
\boxed{
\omega=\omega_0+\alpha t
},
\]

\[
\boxed{
\theta
=
\theta_0+\omega_0t+\frac12\alpha t^2
},
\]

\[
\boxed{
\omega^2
=
\omega_0^2+2\alpha\Delta\theta
}.
\]

Finalmente, el movimiento circular puede describirse completamente mediante
componentes rectangulares:

\[
\boxed{
\vect{r}
=
R\cos\theta\,\ihat
+
R\sen\theta\,\jhat
},
\]

\[
\boxed{
\vect{v}
=
-R\omega\sen\theta\,\ihat
+
R\omega\cos\theta\,\jhat
},
\]

\[
\boxed{
\vect{a}
=
\left(
-R\alpha\sen\theta-R\omega^2\cos\theta
\right)\ihat
+
\left(
R\alpha\cos\theta-R\omega^2\sen\theta
\right)\jhat
}.
\]

Estas relaciones permiten describir la posición, velocidad y aceleración de una
partícula que recorre una circunferencia, conservando la coherencia con la
cinemática vectorial estudiada anteriormente.

\end{resumen}